\section{机构配置解释与研究边界}
\subsection{应用条件}
Global RRP 提供了风险预算参考、收益目标和资本约束相衔接的资产端配置框架，可作为机构开展多资产评估的研究输入。实际应用仍须结合负债期限、现金支付安排和账户约束；资产端历史收益不能直接证明负债覆盖能力。协方差波动也不等同于信用、流动性或监管资本风险，相关要求需要独立数据和适用规则支持。

\subsection{实施限制}
本文按统一费率扣减交易成本，未使用成交回报、盘口或基金申赎资料估计执行能力。实际容量取决于交易金额、市场深度和价格冲击\cite{almgren2001optimal}。周度调仓还可能在收益信号反转或相关性变化时增加损失，因此低历史波动不能替代成交验证和压力评估。逐期权重调整的解释应结合输入变化与目标项贡献，不能仅依据持仓图作因果归因。

\subsection{后续验证}
当前规格经过历史研究选择；逐期输入严格早于决策日，并不消除规格选择偏差。后续应冻结估计、年度校准、资产池复核与计费规则，在新增观察上连续记录持仓、收益和求解失败。若调整规则，应保留生效日期并区分前后证据。本文结论限定为既定样本下的可复现配置结果，不构成未来收益或直接部署的保证。
